\begin{trapbox}
	\begin{description}[style=nextline, leftmargin=2.8em, labelsep=0.5em, itemsep=0.35em]
		\item[\textbf{\textcolor{CTrap}{易错点一（分母开方遗漏）}}] 在计算距离 $d$ 时，常将分母误写为 $A^2+B^2$，遗漏二次根号，导致距离计算错误.
		\item[\textbf{\textcolor{CTrap}{正确理解（公式完整记忆）：}}] 分母是 $\sqrt{A^2+B^2}$，不能省略根号.
		\item[\textbf{\textcolor{CTrap}{错因分析（与斜率混淆）：}}] 学生可能将点到直线的距离公式与直线斜率 $k=-\dfrac{A}{B}$ 混淆，或与两点间距离公式混淆，只记得平方和而忘记开方.
	\end{description}

	\medskip
	\begin{description}[style=nextline, leftmargin=2.8em, labelsep=0.5em, itemsep=0.35em]
		\item[\textbf{\textcolor{CTrap}{易错点二（不等号方向混淆）}}] 判断位置关系时，误认为 $d>r$ 表示相交，$d<r$ 表示相离.
		\item[\textbf{\textcolor{CTrap}{正确理解（大小与位置对应）：}}] 距离越小直线越靠近圆心，所以 $d<r$ 才穿过圆（相交），$d>r$ 则远离（相离）.
		\item[\textbf{\textcolor{CTrap}{错因分析（直觉错误）：}}] 学生可能认为直线“超出”半径才是相交，事实上相交意味着直线伸入圆内，距离小于半径才满足；建议通过画图或口诀“小则交，等则切，大则离”加强记忆.
	\end{description}

	\medskip
	\begin{description}[style=nextline, leftmargin=2.8em, labelsep=0.5em, itemsep=0.35em]
		\item[\textbf{\textcolor{CTrap}{易错点三（忽略竖直直线）}}] 当直线方程为 $x=x_1$（即 $A=1,B=0$）时，部分学生忘记套用距离公式，或错误使用斜截式.
		\item[\textbf{\textcolor{CTrap}{正确理解（统一使用一般式）：}}] 将直线化为一般式 $Ax+By+C=0$，无论 $B=0$ 还是 $A=0$，距离公式 $d=\dfrac{|Ax_0+By_0+C|}{\sqrt{A^2+B^2}}$ 均适用，此时 $\sqrt{A^2+B^2}=|A|$ 或 $|B|$，结果正确.
		\item[\textbf{\textcolor{CTrap}{错因分析（思维定势）：}}] 学生习惯将直线写成 $y=kx+b$，遇到竖直直线时不知如何入手；应养成使用一般式的习惯，把 $x=2$ 转化为 $x-2=0$ 再代入公式.
	\end{description}
\end{trapbox}
